
\chapter{Conclusioni}

% Discutere di possibili miglioramenti futuri --> test su drone reale --> inserimento di altri algoritmi presentati come quelli per correggere DVL o calibrazione iniziale sensori


In questo paragrafo ci sarà la rappresentazione grafica di tutte le traiettorie di validazione e test per la v4.1.0 che verrà indicata come quella più prestazionale alla fine dell'analisi condotta nel precedente capitolo. Qui vorrei anche evidenziare l'aumento dell'errore in pochi secondi per le varie traiettorie come osservato durante le presentazioni dei risultati.

Ci sarà poi un riassunto di quanto fatto, eventualmente un confronto numerico del miglioramento ottenuto tra la v4.1.0 e la 1.1.0 (struttura base) o direttamente con i risultati del paper originario per le traiettorie su cui è possibile farlo. 

Infine indicherò alcune possibili strade future per il miglioramento, a partire dalla raccolta dati in ambiente reale all'uso di algoritmi IA per calibrazione dei sensori o miglioramento dati DVL in concomitanza con NavNet (con alcune criticità legate all'inferenza relativa da valutare) e potenzialmente proporre un approccio alternativo in cui si usa anche un EKF come fosse una misurazione da fornire a NavNet. Ovviamente anche ulteriori test sulla rete modificando anche elementi di carattere strutturale che potenzialmente può portare ad un miglioramento. 